\section{敏感性、稳健性与失效区域}

\subsection{风险分位数}

\begin{longtable}[]{@{}rrr@{}}
\toprule
风险分位 q & Q2总费用/元 & Q3总费用/元 \\
\midrule
\endhead
0.7 & 15,148,998.173 & 14,818,738.874 \\
0.8 & \textbf{14,916,093.656} & \textbf{14,679,141.979} \\
0.9 & 14,952,447.145 & 14,776,266.569 \\
\bottomrule
\end{longtable}

0.8 在两类模型中均为三组候选最低值，说明风险储备过低会增加 5 倍紧急购电，过高又会增加正常购电和储能占用，两者之间存在稳定的成本平衡。

\begin{figure}[H]
\centering
\includegraphics{../figures_reaslab/fig05_quantile_sensitivity.pdf}
\caption{风险分位数敏感性}
\end{figure}

\subsection{数值可行性与稳健性}

所有正式模型均通过以下确认性检查：SOC 位于 1200--10800 kWh；单个 10 min 充/放电不超过 833.333 kWh；同时充放电次数为 0；跨日 SOC 连续误差为 0；加入紧急购电后，实际供需平衡最小供需松弛量仅为 $10^{-13}$ 量级的浮点误差。问题四因果策略亦独立执行同样的可行性检验。

\subsection{可能失效的区域}

\begin{enumerate}
\item \textbf{星期周期突变。} 节假日、极端天气或社区活动会削弱 D-7 负载和价格预测。
\item \textbf{连续阴雨与分布漂移。} 28 d 残差分位假设近期误差分布具有一定稳定性；若光伏误差突然发生状态切换，风险裕度会滞后。
\item \textbf{SOC 长时间触边。} 极端价格或持续光伏富余可能使容量约束成为主导，此时结果对储能容量和效率更敏感。
\item \textbf{储能退化未计。} 当前目标没有循环寿命/退化成本，可能高估频繁套利价值；若给出退化参数，可在线性目标中增加吞吐量成本。
\item \textbf{终端 SOC 价值。} Q2--Q4 采用日内优化并跨日传递 SOC，没有强制每日终端等式；若运营方要求为次日留出固定备用，应增加终端 SOC 价值或滚动 24--48 h 预测。
\item \textbf{结算解释。} Q3 下调价格规则存在文字歧义；本文以敏感性实验约束结论强度，不把绝对费用推广为唯一答案。
\end{enumerate}

\section{模型评价与创新点}

本文没有把“用了复杂算法”作为创新，而是把创新集中在题目真实困难上：

\begin{itemize}
\item 通过紧急购电与正常购电的非对称成本，把 5 倍惩罚转化为具有经济含义的 \textbf{0.8 分位风险储备}，并用敏感性实验验证；
\item 对问题三不是机械使用全部外部预报，而是依据过去 28 d 的 MAE 在“外部预报/历史预报”间滚动择优，使 0:00 与 18:00 的低价值外部信息可以自动被拒绝；
\item 对问题四主动识别并修复“未来真实价格泄露”，构造\textbf{因果策略 + 完全信息下界（oracle）}双层评价，从而同时回答可实施成本和信息价值；
\item 全部主要结论均由对应模型、数值实验与约束检验共同支持，并在结论中按证据强度控制表述范围。
\end{itemize}

模型的主要优点是结构统一、参数含义清晰、全部优化均为 LP、无随机种子、能复现并审计；不足是负载/PV/价格预测仍较简单，且未加入储能退化、线路损耗与需求响应。当题目提供更丰富的气象或价格驱动变量时，可在不改变调度层的前提下替换预测层。

\section{结论}

本文从确定性典型日出发，逐步加入净负荷不确定性、多时刻光伏预报和波动电价，形成统一的微网购电--储能滚动优化框架。问题一表明储能可显著利用价差与光伏错峰，典型日电费下降 26.90\%；问题二通过 0.8 分位风险储备在常规购电和 5 倍紧急购电之间取得平衡；问题三证明 6:00、12:00 的新增光伏预报值得用于调整，并在全年尺度将紧急购电量降低 28.98\%、总费用降低 1.59\%；问题四在严格因果价格预测下仍保持约 1.69\% 的滚动调整收益，而且距完全价格信息 oracle 仅约 1.1\%。这些结论均限定在当前数据、效率口径和结算解释下，不作超出证据范围的普适性推断。
